\documentclass[11pt,a4paper]{article}
\usepackage[margin=2.2cm]{geometry}
\usepackage{amsmath,amssymb,xcolor,enumitem,booktabs,graphicx,tabularx}
\usepackage[normalem]{ulem}
\setlength{\parskip}{0.6em}\setlength{\parindent}{0pt}
\definecolor{newc}{HTML}{1f6f43}
\newcommand{\new}[1]{\textcolor{newc}{#1}}
\newcommand{\old}[1]{\textcolor{gray}{\sout{#1}}}
\newcommand{\note}[1]{{\small\itshape [#1]}}
\newcommand{\lab}[1]{\textsf{\footnotesize #1}}
\renewcommand{\ref}[1]{\lab{#1}}
\renewcommand{\eqref}[1]{(\lab{#1})}
\newcommand{\nref}[1]{the (\lab{#1})}
\newcommand{\citep}[1]{[\lab{#1}]}
\newcommand{\citet}[1]{\lab{#1}}
\title{Proposed change, ``How the controllers learn''}
\date{24 November 2026, version 3}
\begin{document}\sloppy
\maketitle

\section*{How to read this}
This answers what is left of \texttt{method\_training\_comments.md} for
\lab{sec:method:training} in \texttt{4.Implementation.tex}. A1, A2, C4, D4
and E were settled by [A109] and [A114]. No value and no behaviour of the
training changes.

Changed from version 1, as agreed on 24 November:
\begin{itemize}[nosep]
  \item The training is told as it was done. Every 16-agent policy trains in
    two stages, a short run and then a long one that starts from its
    weights. The Method gives only reasons from the design. How the runs
    went, and why training became longer, is in Appendix~J.
  \item The credit is introduced with PRIMAL, then ``here the group is
    instead''. ``The groups used here are the thesis's own'' and PRIMAL's
    ``no clear change'' are gone.
  \item The crowding paragraph is cut. The Problem Formulation already says
    it (\lab{sec:pf:controllers}, ``a decentralised agent is not charged for
    the delay it causes others'').
  \item The Discussion gets one new point, training from random weights on
    long episodes.
  \item Version 3: Appendix~J tells how the second stage came about, as the
    author put it. The first runs fell behind, the episodes were too short,
    and the policy had not learned enough in the budget.
  \item The Colab restart of five runs after 144\,000 steps is left out of the
    thesis. The weights carried over and each run still trained
    2\,000\,000 steps, so ``every run is trained the same way'' holds. It is
    recorded in the GAPS entry.
\end{itemize}

Black text stays, \new{green} is new and \old{struck} text goes.
\note{Notes in brackets} say why. Nothing below is in the thesis yet.

\section*{Part 1. What was checked}
\begin{itemize}[nosep]
  \item \texttt{training/env.py}: \texttt{team\_mean} replaces each reward by
    the mean over the group, for the centralised and section-based
    controllers only. The decentralised controller keeps its own reward.
    \texttt{\_start\_timestep} fixes the groups at the start of the
    timestep, before anyone moves.
  \item \texttt{models/base\_model.py}: the value is its own network
    (Linear, ReLU, Linear) over $z_i$ alone.
  \item \texttt{training/config.py}: three seed sets (training, validation,
    test).
  \item \texttt{SESSION\_CONTEXT.md}, 19 November: every 16-agent long run
    starts from the checkpoint its 500\,000-step run was tested with
    (\texttt{train --init-checkpoint}, \texttt{feat/warm-start-a93}), then
    2\,000\,000 steps on episodes of 16\,000, 3 runners, so 48\,000 steps per
    iteration. The checkpoint is chosen afterwards on the 5 validation seeds.
  \item The pilot (decentralised with messages, fixed storage): on the 20
    test seeds the long run reached 0.156 tasks per step against 0.091,
    served share 0.56 against 0.37. On runs of 100\,000 steps the earlier
    policy broke down on all 3 seeds within 10\,000 to 20\,000 steps.
  \item An evaluation run lasts 10\,000 steps (\lab{sec:exp:setup}).
  \item PRIMAL (\texttt{papers/1809.03531v3.pdf}): ``apply joint rewards to
    agents in an attempt to help them realize the benefit of taking personal
    sacrifices to benefit the team. We briefly tried to assign joint rewards
    to agents within the same FOV.'' It does not say mean or sum.
\end{itemize}

\section*{Part 2. The section as it would read}

\subsection*{How the controllers learn}

Every learned controller is trained with PPO \citep{schulman2017ppo}\old{,
with one policy shared by all agents, as RLlib supports
\citep{rllibmultiagent}}\new{. All agents of a controller share one policy,
the common recipe of Section~\ref{sec:bg:marl}}. \old{Seeds are split between
training and evaluation before any run starts.} Each controller is trained
separately with the same budget of environment steps, so a difference between
them is not a difference in training effort. One environment step is one
timestep of the warehouse, however many agents choose within it. The
centralised controller therefore gets no more experience for running its
network once per agent. \old{Its extra cost shows in runtime and training time
instead} \new{It pays in time instead}
(Section~\ref{sec:impl:cost}). The PPO values, the budget and how a checkpoint
is chosen are in Appendix~\ref{app:training}. \note{B2. The idea is the
Background's, RLlib is the tool and stays in the Simulator section. D3. The
seed split moves to the next paragraph, and runtime is told once, in
``Choosing together''.}

\new{Each seed is its own warehouse. Before any run starts the seeds are split
three ways. Training seeds are played in training, validation seeds choose the
checkpoint and test seeds give the results.} Training plays a new warehouse in
every episode\old{, 250 in the budget of 500\,000 steps in episodes of 2000.
Each controller sees as many layouts as the budget allows, and}\new{, so}
\old{it}\new{a policy} cannot learn one map by heart.
\old{The 16-agent runs were then trained for 2\,000\,000 more steps on
episodes of 16\,000 (Appendix~\ref{app:training}).}
\new{Training has two stages, the same for every controller and storage rule.
A first run plays 500\,000 steps on episodes of 2000. A second run starts from
the weights of the first and plays 2\,000\,000 steps on episodes of 16\,000.
Starting there saves training time. The episodes of the second run are longer
than an evaluation run, 10\,000 steps, so a fleet that falls behind over an
evaluation run also falls behind in training and in the validation that
chooses the checkpoint. How the two stages were chosen is in
Appendix~\ref{app:training}.} \note{C3, the three names of the code. The
training as it was done, with reasons from the design only. ``Same budget''
now means both stages, 2\,500\,000 steps in all.}

Each agent's choice is one training sample with its own value estimate. It
is credited with the mean reward of \new{its \emph{credit group}}\old{the
agents that chose together with it at that timestep}. \new{A joint reward of
this kind lets an agent gain from helping the others.
\citet{sartoretti2019primal} tried one for the agents within one field of
view. Here the group is instead the agents that chose together with it at
that timestep.} These are the whole fleet for the centralised controller, the
agents in the same zone for the section-based controller, and the agent alone
for the decentralised controller. A zone's agents change as agents cross its
border, and the \old{mean is taken over the agents in the zone at
$t$}\new{group is fixed by where the agents stand at the start of $t$, before
anyone moves}. \new{The credit is a mean rather than a sum, so it stays on the
scale of one agent's reward whatever the size of the group, and all three
controllers can train with the same PPO values.} Formally, with
$\mathcal{A}_i(t)$ the \old{agents that choose together with}\new{credit group
of} $a_i$ at $t$, the \emph{credit} of agent $a_i$ is
\begin{equation*}
  \bar R_i(t) = \frac{1}{|\mathcal{A}_i(t)|} \sum_{a_j \in \mathcal{A}_i(t)} R_j(t) .
  \tag{\lab{eq:credit}}
\end{equation*}
\note{B1, the joint reward is borrowed, the group and the mean are named
without a label. C2, the symbol stays and the word is the code's. What the
credit means for each controller is in the Problem Formulation and is not
told again.}

The value estimate comes from a value head $V_\theta$, a small fully
connected network with ReLU activations over the same vector $z_i$ as the
policy. So the critic\old{, the part of the learner that estimates how good a
situation is,} sees the same window as its own policy\old{. There are two
reasons. What}\new{, because what} a controller sees is set by the agents it
decides for, so that limit should hold when it learns as well as when it
acts. \old{It also keeps the difference in crowding between the controllers
visible. A decentralised learner gets no signal about the crowding it causes
for agents it cannot see, and that is the effect the comparison is meant to
show. A critic that sees the whole warehouse during training only would
partly teach that cost and hide the difference, so it is not used.} \new{A
critic that sees the whole warehouse during training, as in centralised
training with decentralised execution \citep{lowe2017maddpg}, is therefore
not used.} \note{D2, the Background explains the critic. A3, one reason, the
crowding is the credit's and the PF has it. B3, the rejected critic named.}

Each storage rule gets its own policy, trained under the same rule it is
tested with, because a warehouse would train its fleet under the rule it
runs. The gain then compares two whole systems, storage rule and policy
together, and part of it can come from training. \old{Training under demand
and congestion-aware storage is also harder, because these rules change which
shelves the picks come from, through \nref{eq:coupling}, while the policy is
learning (Section~\ref{sec:pf:coupling}).} \note{B4. No result shows it, so
it goes.}

\section*{Part 3. Other places that change}

\textbf{Appendix J, training budget (\lab{app:training}).} \note{The outcomes
that set the two stages belong here, in the experiment paper, not in the
Method.}

\old{Every learned controller is trained for 500\,000 environment steps on
episodes of 2000 steps, at every fleet. The budget was not read off learning
curves. The 16-agent runs were then trained again, from a warm start at the
checkpoint each run had been tested with, for 2\,000\,000 more steps on
episodes of 16\,000 steps. That is 125 episodes over the 250 training
warehouses.}

\new{The first stage trains every learned controller for 500\,000 environment
steps on episodes of 2000, at every fleet. That is about one episode on each
of the 250 training warehouses. The budget was not read off learning curves.}

\new{At 16 agents the first-stage policies fell behind in the tests. A pilot
on the decentralised controller with messages under fixed storage looked for
the cause, and found two. The episodes were too short. Played for 100\,000
steps on 3 seeds, the first-stage policy broke down within 10\,000 to 20\,000
steps, later than a training episode of 2000 ever reaches. The budget was too
small. A second run from its weights, on episodes of 16\,000, kept learning.
Its throughput in training rose from 0.079 to 0.184 tasks per step over the
first 18 of its 42 iterations. On the 20 test seeds it raised the throughput
from 0.091 to 0.156 tasks per step and the served share from 0.37 to 0.56.
The pilot changed the budget and the episode length together, so it cannot
say which of the two helped more.}

\new{So every 16-agent controller and storage rule gets the same second stage.
It starts from the checkpoint its first run was tested with and plays
2\,000\,000 steps on episodes of 16\,000, 125 episodes over the 250 training
warehouses. The second stage starts from the first rather than from random
weights to save time. One second-stage run is about 42 iterations of 48\,000
steps. The first-stage runs are kept as earlier runs.} \note{Pilot numbers
from \texttt{SESSION\_CONTEXT.md}. The episode-length plot (2000, 8000,
16\,000) and the list of every run come with the Appendix~L ``earlier runs''
rewrite.}

\textbf{Experiments, ``What could have been done differently''.}

Longer training episodes were the first change, and they helped.
\old{Four}\new{Six} more could bring the learned controllers closer to Token
Passing. \note{The list already had five items, and gets one more.}
\begin{itemize}[nosep]
  \item \emph{More training seeds.} Each run had one seed. \ldots
  \item \new{\emph{Training from random weights.} Every 16-agent policy
    started its long run from its first run, so a weak first run could hold
    its long run back. Runs from random weights on episodes of 16\,000 would
    show how much the start matters.}
  \item \emph{A way out of locks.} \ldots A shared \old{reward}\new{credit}, or
    a rule that makes one of two locked agents give way, would let them back
    out.
\end{itemize}

\textbf{``When moves still clash'' (\lab{sec:method:resolution}).} \note{D1,
the credit is told once, in the training section.}

\ldots A decentralised agent\old{ is credited with its own reward only}\new{'s
credit is its own reward only (Section~\ref{sec:method:training})}, so it
gains nothing when it backs out for another agent. The section-based and
centralised controllers share the \old{reward of the agents that choose
together}\new{credit of their group}, so they have a reason to learn it.
\ldots

\textbf{Notation table, row \lab{eq:credit}.}

\old{The agents that choose together with $a_i$ at $t$}\new{The credit group
of $a_i$ at $t$, the agents that choose together with it}, and the mean of
their rewards, which trains $a_i$'s choice.

\textbf{\texttt{references.bib}.} \texttt{author = \{\{Ray Team\}\}} in both
RLlib entries, so the citation prints as ``Ray Team, 2024''.

\section*{Questions}
\begin{enumerate}[nosep]
  \item Appendix~J now gives the pilot numbers as the reason for the second
    stage. Enough, or should the episode-length plot come in this change
    too?
  \item The pilot's test numbers (0.156, 0.56) come from a checkpoint chosen
    on 2 validation seeds by mistake (\texttt{*\_2seeds}). Was it chosen
    again on 5 seeds (\texttt{eval-comm-long-v5})? If so, which numbers go
    here?
  \item Some second-stage runs are not evaluated yet (\texttt{TODO [A110]} in
    the Experiments). The Method describes the training as done for all of
    them. Fine to apply now?
\end{enumerate}

\end{document}
